\section{Question 4: Machine Translation}

\subsection{Dataset Description}

Question 4 uses the English-to-German Multi30k translation benchmark \cite{elliott2016multi30k}. In this project, the dataset is accessed through its Hugging Face mirror and treated as a short-caption translation task. That makes it a useful benchmark for testing whether a model can produce fluent target-language sentences while preserving compact visual descriptions from the source caption. The current Q4 evidence now includes two finished systems evaluated on the same matched 200-example validation and 200-example test splits: a pretrained transformer reference and a custom seq2seq-plus-attention baseline.

The comparison is direct at evaluation time but not perfectly budget-matched at training time. The transformer run keeps a 100-example train slice only for artifact-contract consistency and does not perform task-specific fine-tuning, whereas the seq2seq model is trained from scratch on 2000 capped training examples. This distinction matters for interpretation, so the section should be read as a practical pretrained-versus-classical comparison rather than a matched from-scratch study.

\subsection{Preprocessing and Model}

The stronger system uses a pretrained Opus-MT English-to-German checkpoint from Helsinki-NLP \cite{tiedemann2020opusmt}. It is built on a transformer-style translation stack \cite{vaswani2017attention}. Source and target casing are preserved, inputs and decoded outputs are capped at 96 tokens, and generation uses beam search with 4 beams and length penalty 1.0. Instead of fine-tuning on Multi30k, this pipeline evaluates the checkpoint directly as a strong practical reference.

The classical comparator is a compact GRU encoder-decoder with additive attention \cite{bahdanau2015nmt}. It uses 128-dimensional token embeddings and hidden states, a single recurrent layer, teacher forcing during training, and a maximum decoded length of 32 tokens. Unlike the transformer reference, this seq2seq model is trained from scratch on a lowercased 2000-example training split. It must therefore learn translation behavior directly from the capped Multi30k data.

Evaluation uses BLEU, ChrF, METEOR, and BERTScore on both validation and test. BLEU gives a familiar n-gram precision view of translation quality, ChrF is more tolerant of inflectional and morphological variation, METEOR accounts for word order and partial matching, and BERTScore captures semantic similarity through contextual embeddings. Together, the four metrics provide a credible multi-faceted evaluation for the current Q4 comparison.

\subsection{Results}

Table~\ref{tab:q4_overall_results} summarizes the finished comparison. The pretrained transformer reaches validation/test BLEU 0.3749/0.3711, ChrF 0.6421/0.6398, METEOR 0.6516/0.6467, and BERTScore 0.8993/0.8910. The seq2seq baseline reaches BLEU 0.1110/0.1364, ChrF 0.3293/0.3448, METEOR 0.3290/0.3544, and BERTScore 0.7269/0.7351. On the shared evaluation the transformer therefore leads by roughly 0.26 BLEU, 0.29 ChrF, 0.29 METEOR, and 0.16 BERTScore on the test split.

\begin{table}[H]
\centering
\small
\begin{tabularx}{\linewidth}{>{\raggedright\arraybackslash}X c c c c c c c c c}
\toprule
Model & Train & Val BLEU & Test BLEU & Val ChrF & Test ChrF & Val METEOR & Test METEOR & Val BERTScore & Test BERTScore \\
\midrule
Opus-MT Transformer & 2000 & \textbf{0.3749} & \textbf{0.3711} & \textbf{0.6421} & \textbf{0.6398} & \textbf{0.6516} & \textbf{0.6467} & \textbf{0.8993} & \textbf{0.8910} \\
Seq2Seq + Attn      & 2000 & 0.1110 & 0.1364 & 0.3293 & 0.3448 & 0.3290 & 0.3544 & 0.7269 & 0.7351 \\
\bottomrule
\end{tabularx}
\caption{Q4 Multi30k EN$\to$DE comparison on the matched 200-example val/test split. Both models trained or evaluated on 2000 training examples. Transformer (Opus-MT) is a pretrained reference; Seq2Seq is trained from scratch.}
\label{tab:q4_overall_results}
\end{table}

Figure~\ref{fig:q4_test_metric_comparison} makes the same comparison easier to scan visually. The transformer leads the classical seq2seq model on all four test metrics, and the gap is large enough that the model ranking is unambiguous even on the capped evaluation split.

\begin{figure}[H]
	\centering
	\includegraphics[width=0.88\textwidth]{figures/q4/test_metric_comparison.png}
	\caption{Q4 test-set metric comparison for the Seq2Seq baseline and the pretrained Opus-MT transformer.}
	\label{fig:q4_test_metric_comparison}
\end{figure}

The size of this gap still makes the main takeaway straightforward: under the current capped setup, the from-scratch seq2seq baseline improves materially with more training budget but is still not competitive with a strong pretrained translation checkpoint. At the same time, the comparison should not be oversold as a perfectly aligned training-budget experiment, because the transformer benefits from large-scale pretraining rather than the capped Multi30k train split.

\subsection{Qualitative Analysis}

The qualitative examples reinforce the metric gap. On test example \texttt{test\_63}, the transformer exactly matches the reference translation for \textit{Two men playing guitar in front of a large audience.}, which is the ideal behavior for a short caption translation task. Even its weaker exported example, \texttt{test\_8}, remains understandable: it rewrites the running-man caption as a longer relative-clause construction that is less concise and less reference-like than the gold translation, but it still preserves most of the core meaning.

The seq2seq outputs show why the classical baseline still remains far behind despite the stronger run. Its best exported test example, \texttt{test\_48}, captures much of the source meaning in \textit{ein älterer mann spielt saxophon. ein videospiel.}, but it still introduces an incorrect content word and awkward sentence structure. Its weakest exported example, \texttt{test\_94}, drifts badly into \textit{drei chefköche, darunter auf dem rücken zu zu}, which is grammatical only in fragments and no longer faithful to the source sports caption. In other words, the seq2seq model now recovers more partial meaning than the weaker earlier run, but it still fails to sustain fluent, source-faithful translation consistently.

\subsection{Discussion}

The Q4 evidence supports a clear conclusion. The pretrained transformer is the strongest translation system by a wide margin and should be treated as the primary Q4 result, while the seq2seq model serves as the classical reference point that shows how difficult translation remains for small from-scratch models on limited data.

The remaining caveat is experimental alignment. Because the transformer is evaluated as a pretrained system rather than fine-tuned on Multi30k, this section should be interpreted as a pretrained-reference comparison, not a strictly matched training-budget contest. The current evidence still answers the assignment question clearly: pretraining dominates the classical seq2seq baseline under the evaluated setup. A follow-up fairness study should fine-tune the transformer on the same train split and report parameter and compute budgets side-by-side.